\section{10 001st Prime --- Problem 7}
\label{prob:7}
\noindent\textbf{Source:} \url{https://github.com/KO0I/projectEuler/blob/reorganize/solutions-by-language/solutions/c/10_001st_prime_7.c}

\noindent\textbf{Assessment:} Does not reliably solve the named problem.

\subsection*{Mathematical goal}
The aim is to identify \(p_{10001}\), the prime in position 10001 of the ordered sequence \(2,3,5,7,\ldots\).

\subsection*{How the C code works}
The program seeds a table with 2 and 3, examines successive odd candidates using trial division up to a rounded square root, and writes accepted candidates to a file named \verb|primes|. Comments discuss using the table for other Euler problems as well.

\subsection*{Correctness and limitations}
The hard-coded \verb|countDone=78497| is unrelated to the 10001st prime, and the desired table entry is never printed. Both the generation loop and the output loop access index \verb|countDone| in an array of that length, which is out of bounds. Candidate-reset behavior during trial division also needs correction. The expected prime is 104743; this file does not establish it.

